\chapter{Presentation and Viva Preparation}
\label{ch:viva}

\section{A concise project introduction}
\begin{quote}
We implemented a six-axis Yaskawa GP7 pick-and-place system in ROS 2 Jazzy
and Gazebo Harmonic. The runtime parses the actual URDF for forward
kinematics, uses joint-limit-bounded numerical least squares for inverse
kinematics, samples a quintic-scaled Cartesian route, and sends it through a
ROS 2 joint trajectory controller. Gazebo pose and grasp feedback are checked
throughout the task. Our fixed Box A position is reachable under the
lecturer's subsequent verbal clarification; the original sheet prints an
unreachable $(1,1,0)$ destination. The successful simulated carry uses an
idealized detachable fixed joint, so we do not claim physical grasp
performance.
\end{quote}

\section{Suggested presentation sequence}
\begin{enumerate}
 \item State the assignment and clarify the original and verbally approved
 Box A coordinate distinction.
 \item Show the robot, identify the six arm joints, and point out the grasp
 frame and adaptive gripper.
 \item Explain the standard-DH frame construction and table, especially the
 $q_2-\pi/2$ offset and signed $d_4=-0.440$ m.
 \item Demonstrate one home/nonzero FK transform and say the executable FK
 reads URDF transforms.
 \item Explain numerical IK residual, bounds, warm-starting, restart seeds,
 and why it is not an analytical GP7 solution.
 \item Show the quintic time law, waypoint sequence, controller action and
 feedback architecture.
 \item Play the Gazebo demonstration. Label far/near mode and seed and narrate
 approach, grasp, carry, release, and retreat.
 \item Present recorded test metrics and immediately qualify the fixed-joint
 grasp, mimic warning, estimated inertias, and absence of physical validation.
\end{enumerate}

\section{Demonstration narration}
Before pressing play, show the fixed Box A at $(0.55,0.35,0)$ and state that
Box B's pickup is the randomized test pose. For near-mode seed 7, say that it
produces the repeatable pickup $(0.532383276,-0.217932066,0.025)$ m; the fixed
destination remains unchanged. Narrate the downward approach and descent,
position-checked close, plugin attachment, lift and upright transfer, lowering
to 1 mm above the center target, release, retreat, and settling. Finish by
showing \code{PICK AND PLACE VERIFIED} and the final error. Explain that the
plugin fixed joint intentionally retains the cube and does not demonstrate
force-based grasping. The report draft in \file{docs/report.md} contains a
step-by-step recording checklist; a headless validation log is not a substitute
for the assignment's screen-recorded video.

\section{Who should understand what}
The task sheet allows at most three group members. The guide does not assign
names or roles. For preparation, all members should be able to explain the
project scope, clarify Box A/Box B, interpret the DH table, and describe the
grasp limitation. A group may distribute deeper preparation as follows, then
fill the placeholders with actual members:
\begin{itemize}
 \item \textbf{Member [name/index]}: URDF, frame assignment, DH derivation,
 FK verifier, and physical meaning of signed parameters.
 \item \textbf{Member [name/index]}: numerical IK, trajectory sampling,
 controller interfaces, and code-to-mathematics mapping.
 \item \textbf{Member [name/index]}: launch architecture, Gazebo world,
 gripper/plugin, test evidence, setup, and limitations.
\end{itemize}
This is a study suggestion only; no group membership is inferred.

\section{Likely viva questions and grounded answers}
\subsection{Foundation}
\begin{enumerate}
 \item \textbf{How many degrees of freedom does the modeled arm have?} Six
 actuated revolute arm joints, named \code{joint_1_s} through
 \code{joint_6_t}. The gripper driver is an additional actuated joint but
 is not counted as one of the six arm axes.
 \item \textbf{What does FK calculate?} The pose of \code{grasp_frame} for
 a known six-element joint vector. The runtime follows fixed URDF origins
 and joint-axis rotations.
 \item \textbf{What is the difference between a ROS topic and action here?}
 Joint states and stop are topics; a trajectory is a long-running action
 with goal acceptance, result, and cancellation.
 \item \textbf{What is the DH convention?} Standard DH:
 $R_z(\theta_i)T_z(d_i)T_x(a_i)R_x(\alpha_i)$. Our frames follow that order.
 \item \textbf{Why is the object kept upright?} Cartesian segments hold
 $R_{down}=\mathrm{diag}(1,-1,-1)$, and box-pose feedback is compared with
 its initial rotation.
\end{enumerate}

\subsection{Intermediate}
\begin{enumerate}
 \item \textbf{Why is $d_4$ negative?} Frame $z_3$ points in $-X$, but the
 wrist-center displacement from frame 3 is in $+X$. DH parameters are signed
 coordinate displacements, hence $d_4=-0.440$ m.
 \item \textbf{Why does $q_2$ have an offset?} The chosen DH frame's zero
 orientation differs from the URDF shoulder zero by $\pi/2$. The mapping is
 $\theta_2=q_2-\pi/2$; no joint sign is reversed.
 \item \textbf{Does the task execute the DH FK?} No. Runtime FK parses URDF;
 the independent documentation script separately evaluates DH and compares
 transforms over 104 configurations.
 \item \textbf{How does the IK solver handle joint limits?} It reads lower
 and upper bounds from URDF, clips each initial seed slightly inside the
 bounds, and calls SciPy bounded least squares.
 \item \textbf{Why seed each IK point with the preceding solution?} Nearby
 seeds favor a continuous branch; the path also rejects adjacent changes
 above 0.15 rad. Neither guarantees global branch optimality.
 \item \textbf{Why use quintic blending?} The scalar polynomial has zero
 endpoint first and second derivatives for rest-to-rest moves. Actual
 Cartesian segments then solve nonlinear IK, so sampled joint acceleration
 is not analytically guaranteed by that polynomial alone.
 \item \textbf{What does the 0.5 rad/s limit mean?} The planner caps
 sampled velocity by the lesser of 0.5 rad/s and URDF per-joint velocity,
 stretches time, and rescales velocity. It does not enforce acceleration.
\end{enumerate}

\subsection{Advanced}
\begin{enumerate}
 \item \textbf{Is this an analytical spherical-wrist IK?} No. It is a
 bounded residual-based SciPy least-squares solve. The source has no explicit
 geometric Jacobian, wrist decomposition, or enumeration of analytical
 branches.
 \item \textbf{What is questionable about the residual norm?} It combines
 metres and radians without normalization, so a single $10^{-6}$ norm
 threshold is a solver acceptance rule, not a dimensionally uniform physical
 tolerance.
 \item \textbf{What does 104-configuration agreement prove?} That the
 documented DH table, fixed tool transform, and runtime URDF FK are
 numerically consistent for four specific and 100 seeded in-limit cases.
 It does not prove the URDF matches manufacturer CAD or physical hardware.
 \item \textbf{Why does a micrometre-scale position error not imply that
 accuracy on a real robot?} The object is rigidly attached by a simulator
 fixed joint and pose is computed in simulation; there is no contact-force,
 sensor, actuator calibration, or hardware measurement.
 \item \textbf{What does the mimic warning mean?} The selected physics
 engine does not support the configured mimic constraints. The task's
 successful object carry comes from detachable-joint retention, so finger
 coupling is not validated by that result.
 \item \textbf{Is /gp7/stop an emergency stop?} No. It requests action
 cancellation and retains an attached simulation object. It is not a
 safety-rated hardware stop.
 \item \textbf{How does the task avoid an IK failure after motion starts?}
 It calculates the approach and every Cartesian segment before sending the
 first arm command; failures in IK or branch continuity abort pre-motion.
 This is not collision preflight.
 \item \textbf{How was reachability evaluated for far mode?} SLSQP
 maximizes horizontal grasp radius subject to $z=0.025$, fixed downward
 orientation, and URDF joint bounds, starting from several angular seeds.
 It backs off the computed boundary by 5 mm. This is pose-specific and the
 exact boundary is straight-elbow singular.
\end{enumerate}

\section{Common explanations to avoid}
\begin{itemize}
 \item Do not say the original assignment sheet changed; state that the
 lecturer later gave a verbal clarification.
 \item Do not say Box A is randomized. Box A stays fixed; near mode samples
 Box B pickup coordinates.
 \item Do not say the runtime uses the DH matrices. It uses URDF FK; DH is
 an independently verified model for explanation.
 \item Do not call the IK analytic, Jacobian-pseudoinverse, or collision-aware.
 \item Do not claim the fingers physically grasp the cube by friction. The
 simulation uses a detachable fixed joint and mimic support is limited.
 \item Do not call sampled velocity a proven acceleration/jerk profile or
 claim the controller's commanded path is a measured physical path.
 \item Do not present a successful Gazebo node log as the required video.
\end{itemize}
